% =====================================================================
%  CHƯƠNG 1: TỔNG QUAN ĐỀ TÀI
% =====================================================================
\chapter{Tổng quan đề tài}
\label{chap:intro}

\section{Đặt vấn đề}

Thương mại điện tử hiện đại đặt người dùng trước hàng triệu lựa chọn, dẫn tới hiện
tượng \textit{quá tải thông tin} (information overload). \textbf{Hệ thống gợi ý}
(Recommender System) ra đời nhằm tự động lọc và đề xuất cho mỗi người dùng những sản
phẩm phù hợp nhất với sở thích của họ. Đối với nhóm sản phẩm \textbf{Software} (phần
mềm) trên Amazon, bài toán càng quan trọng bởi sản phẩm mang tính chuyên biệt cao,
mô tả kỹ thuật dài và đánh giá người dùng giàu ngữ nghĩa.

Tuy nhiên, các hệ thống gợi ý truyền thống dựa trên \textbf{định danh} (ID-based
collaborative filtering) bộc lộ hai hạn chế cố hữu:

\begin{itemize}[leftmargin=2em, itemsep=2pt]
    \item \textbf{Bỏ qua dữ liệu văn bản ngữ nghĩa:} các mô hình lọc cộng tác cổ
    điển chỉ học từ ma trận tương tác (ai đã mua/đánh giá gì), hoàn toàn không khai
    thác phần mô tả sản phẩm, tính năng hay nội dung review -- vốn chứa thông tin
    rất giá trị về bản chất sản phẩm.
    \item \textbf{Bài toán khởi động lạnh (cold-start):} khi xuất hiện người dùng
    mới hoặc sản phẩm mới chưa có lịch sử tương tác, mô hình dựa trên định danh không
    thể đưa ra dự đoán đáng tin cậy.
\end{itemize}

\section{Giải pháp đề xuất}

Đồ án đề xuất một hướng tiếp cận hiện đại nhằm khắc phục đồng thời hai hạn chế trên,
dựa trên ba trụ cột:

\begin{enumerate}[leftmargin=2em, itemsep=2pt]
    \item \textbf{Mô hình ngôn ngữ (Language Model):} sử dụng một mô hình ngôn ngữ
    dựa trên kiến trúc Transformer để \textit{hiểu} văn bản mô tả sản phẩm và nội
    dung đánh giá, chuyển toàn bộ thông tin văn bản dài thành một \textit{vector số
    cố định} đại diện cho ý nghĩa của sản phẩm. Nhờ vậy, hệ thống nắm bắt được sự
    tương đồng ngữ nghĩa ngay cả với sản phẩm mới.
    \item \textbf{Kiến trúc hai tháp (Two-Tower):} tách riêng việc biểu diễn
    \textit{người dùng} và \textit{sản phẩm} thành hai nhánh độc lập, cùng ánh xạ
    vào một không gian vector chung. Nhánh người dùng học từ \textit{chuỗi hành vi
    theo thời gian}, có tính tới yếu tố thời điểm và mức độ ưa thích.
    \item \textbf{Học tương phản (Contrastive Learning):} huấn luyện mô hình bằng
    cách kéo các sản phẩm mà người dùng thích lại gần vector người dùng và đẩy các
    sản phẩm không liên quan ra xa trong không gian biểu diễn.
\end{enumerate}

\section{Mục tiêu và phạm vi}

\textbf{Mục tiêu} của đồ án là xây dựng một \textit{luồng xử lý hoàn chỉnh}
(end-to-end pipeline) cho hệ thống gợi ý phần mềm, bao gồm:

\begin{itemize}[leftmargin=2em, itemsep=2pt]
    \item Nạp và tiền xử lý dữ liệu lớn bằng \textbf{Apache Spark}, lưu trữ tối ưu
    theo định dạng cột \textbf{Apache Parquet}.
    \item Mã hóa "trải nghiệm" của người dùng thành vector bằng mô hình ngôn ngữ,
    tích hợp trọng số thời gian.
    \item Huấn luyện mô hình hai tháp với hàm mất mát tương phản.
    \item Đánh giá khách quan độ hiệu quả trên \textbf{toàn bộ kho sản phẩm}
    (full-catalog), thay vì chỉ trên một nhóm nhỏ.
\end{itemize}

\textbf{Phạm vi:} Báo cáo tập trung vào danh mục \textit{Software} của bộ dữ liệu
Amazon Reviews 2023, với khoảng $146.000$ người dùng và hơn $17.000$ sản phẩm phần
mềm độc nhất. Bài toán được mô hình hóa theo dạng \textbf{truy hồi} (retrieval):
với mỗi người dùng, hệ thống xếp hạng toàn bộ kho sản phẩm và trả về danh sách
Top-$N$ phù hợp nhất.

\section{Đóng góp của báo cáo}

\begin{itemize}[leftmargin=2em, itemsep=2pt]
    \item Xây dựng pipeline xử lý dữ liệu lớn end-to-end trên PySpark, tạo ra biểu
    diễn văn bản ngữ nghĩa (\texttt{semantic\_text}) thống nhất cho sản phẩm.
    \item Thiết kế thuật toán tổng hợp vector lịch sử người dùng có \textit{trọng số
    thời gian} (recency) kết hợp mức độ ưa thích (rating), tính toán hiệu quả trên
    bộ nhớ.
    \item Cài đặt và huấn luyện mô hình Two-Tower với mô hình ngôn ngữ làm bộ mã hóa
    sản phẩm và hàm mất mát tương phản InfoNCE.
    \item Đề xuất giao thức đánh giá \textbf{full-catalog} loại bỏ thiên lệch đánh
    giá cục bộ, chứng minh mô hình vượt trội nhiều lần so với baseline ngẫu nhiên.
\end{itemize}

\section{Cấu trúc báo cáo}

Báo cáo được tổ chức theo năm chương:

\begin{itemize}[leftmargin=2em, itemsep=2pt]
    \item \textbf{Chương 1 -- Tổng quan đề tài:} đặt vấn đề, giải pháp, mục tiêu và
    phạm vi.
    \item \textbf{Chương 2 -- Cơ sở lý thuyết:} nền tảng dữ liệu lớn, mô hình ngôn
    ngữ, kiến trúc hai tháp và học tương phản.
    \item \textbf{Chương 3 -- Phân tích và thiết kế hệ thống:} kiến trúc tổng thể,
    xử lý dữ liệu, chiến lược chia dữ liệu và thiết kế mô hình.
    \item \textbf{Chương 4 -- Kết quả thực nghiệm:} môi trường, giao thức đánh giá
    và phân tích kết quả.
    \item \textbf{Chương 5 -- Kết luận và hướng phát triển:} tổng kết, hạn chế và
    định hướng tương lai.
\end{itemize}
